\section{Learning Channel Denoising and Activity Detection}
\label{sec:learning}

In this section, we introduce score-based representations of the channel
mean and variance to circumvent the high-dimensional integrals in channel estimation.
We then apply denoising score matching for channel denoising and binary
classification for activity detection.

\subsection{Channel Mean Estimation with a Score-Based Generative Prior}
\label{sec:score-denoising}
\label{sec:channel-learning}

Under the active hypothesis in \eqref{eq:pseudo-awgn},
$\mathbf H\sim q_{\rm tr}$ is observed in independent proper complex
Gaussian noise with covariance $v\I_D$, where $v>0$.
Stacking the real and imaginary parts of $\vect(\mathbf H)$ and
$\vect(\mathbf R)$ into $\mathbf h_{\mathbb R}$ and $\mathbf u$,
respectively, we rewrite the AWGN observation model as
\begin{equation}
 \mathbf u=\mathbf h_{\mathbb R}+\sqrt{v/2}\,\bm\epsilon,
 \qquad \bm\epsilon\sim\mathcal N(\mathbf0,\I_{2D}).
 \label{eq:training-query}
\end{equation}

At fixed $v$, let $q_v^{\mathbb R}$ denote the density of $\mathbf u$.
For the model in \eqref{eq:training-query}, Tweedie's formula
\cite{efron2011,dytso2023} gives the conditional mean as
\begin{equation}
 \mathbf m_{\mathbb R}(\mathbf u,v)
 =\mathbf u+\frac v2\nabla_{\mathbf u}
       \log q_v^{\mathbb R}(\mathbf u),
 \label{eq:tweedie-mean}
\end{equation}
where $\nabla_{\mathbf u}\log q_v^{\mathbb R}$ is the first-order score
function of $q_v^{\mathbb R}$.

For a specific value of $v$, our goal is to learn a first-order score
network $\mathbf s_\theta(\cdot,v):\mathbb R^{2D}\to\mathbb R^{2D}$
parameterized by $\theta$ with the following score matching objective:
\begin{equation*}
 \min_\theta\;
 \E_{\mathbf u\sim q_v^{\mathbb R}}
 \left[\norm{\mathbf s_\theta(\mathbf u,v)
 -\nabla_{\mathbf u}\log q_v^{\mathbb R}(\mathbf u)}^2\right].
\end{equation*}
This objective cannot be directly used as the loss function since
$\nabla_{\mathbf u}\log q_v^{\mathbb R}(\mathbf u)$ is unknown.
The above objective has the same minimizers as the denoising score
matching objective \cite{vincent2011}, whose target is the conditional
Gaussian score $-2(\mathbf u-\mathbf h_{\mathbb R})/v$.

In practice, the score-based MMSE denoiser is required to operate across
a range of noise levels. We therefore combine the denoising score
matching objectives for different values of $v$ with weight $v/(4D)$
to establish a unified objective:
\begin{equation}
 \mathcal L_{\rm DSM}(\theta)=
 \E\left[\frac{v}{4D}
 \norm{\mathbf s_\theta(\mathbf u,v)
 +\frac2v(\mathbf u-\mathbf h_{\mathbb R})}^2\right].
 \label{eq:denoising-score-objective}
\end{equation}
The expectation is taken over the clean channels,
noise variances, and Gaussian noise.

We parameterize the score by a denoiser $\mathbf f_\theta(\mathbf u,v)$
as $\mathbf s_\theta=2(\mathbf f_\theta-\mathbf u)/v$.
Substituting this expression into \eqref{eq:denoising-score-objective}
yields
\begin{equation}
 \mathcal L_{\rm ch}(\theta)=
 \E\left[\frac{1}{Dv}
 \norm{\mathbf f_\theta(\mathbf u,v)-\mathbf h_{\mathbb R}}^2\right].
 \label{eq:first-loss}
\end{equation}
The conditional mean $\mathbf m_{\mathbb R}(\mathbf u,v)$ minimizes
this loss over all estimators of $\mathbf h_{\mathbb R}$ from
$(\mathbf u,v)$. Reshaping the real and imaginary outputs gives the
complex-valued estimate $\mathbf F_\theta(\mathbf R,v)$.

\subsection{Channel Variance Estimation}
\label{sec:conditional-error}

Likewise, the conditional covariance matrix
$\mathbf P_{\mathbb R}=\operatorname{Cov}(\mathbf h_{\mathbb R}\mid\mathbf u,v)$
can be obtained through the second-order Tweedie's formula
\cite{dytso2023} as
\begin{equation}
 \mathbf P_{\mathbb R}
 =\frac v2\I_{2D}+\frac{v^2}{4}
   \nabla_{\mathbf u}^2\log q_v^{\mathbb R}(\mathbf u).
 \label{eq:tweedie-covariance}
\end{equation}

Differentiating \eqref{eq:tweedie-mean} with respect to $\mathbf u$
at fixed $v$ gives the following compact expressions for the covariance
in \eqref{eq:tweedie-covariance} and its normalized trace:
\begin{align}
 \mathbf P_{\mathbb R}&=\frac v2\mathbf J_m,
 \label{eq:mean-derivative-covariance}\\
 \tau^\star&=\frac{v}{2D}\operatorname{div}\mathbf m_{\mathbb R}.
 \label{eq:mean-derivative-variance}
\end{align}
Here $\mathbf J_m=\partial\mathbf m_{\mathbb R}/\partial\mathbf u$
denotes the Jacobian of $\mathbf m_{\mathbb R}$, and $\tau^\star$ is the
average conditional variance per complex coordinate.
This suggests that the average conditional variance can be estimated
by computing the divergence of the trained denoiser $\mathbf f_\theta$
\cite{manor2024}.

The divergence is estimated by a Monte Carlo method \cite{metzler2016},
using $L_z$ independent random sign vectors $\mathbf z_\ell$ with
$\E[\mathbf z_\ell\mathbf z_\ell^{\mathsf T}]=\I_{2D}$.
Central finite differences along $\mathbf w_\ell=\sqrt{v/2}\,\mathbf z_\ell$,
with step $\delta_{\rm FD}>0$, give an estimate of the average conditional variance:
\begin{equation}
 \widehat\tau_\theta(\mathbf R,v)
 =\frac{1}{2\delta_{\rm FD} D L_z}\sum_{\ell=1}^{L_z}
       \mathbf w_\ell^{\mathsf T}\Delta_\ell \mathbf f_\theta,
 \label{eq:derivative-trace}
\end{equation}
where $\Delta_\ell \mathbf f_\theta
 =\mathbf f_\theta(\mathbf u+\delta_{\rm FD}\mathbf w_\ell,v)
 -\mathbf f_\theta(\mathbf u-\delta_{\rm FD}\mathbf w_\ell,v)$.

\subsection{Activity Probability Estimation Based on Density Ratios}
\label{sec:density-ratio}

The density ratio $q_{1,v}/q_{0,v}$ can be estimated from samples with
binary labels $a\in\{0,1\}$, where
$q_{1,v}$ and $q_{0,v}$ are the class-conditional densities at fixed $v$
\cite{menon2016}. For the activity prior $\rho$, Bayes' rule implies
\begin{equation}
 \log\frac{p^\star(\mathbf R,v)}{1-p^\star(\mathbf R,v)}
 =\log\frac{\rho}{1-\rho}
  +\log\frac{q_{1,v}(\mathbf R)}{q_{0,v}(\mathbf R)}.
 \label{eq:activity-log-odds}
\end{equation}

For training, we set $\Pr(a=1)=1/2$.
Samples from the active class $a=1$ are generated as
$\mathbf R=\mathbf H+\mathbf W$, with $\mathbf H\sim q_{\rm tr}$,
whereas samples from the inactive class $a=0$ are generated as
$\mathbf R=\mathbf W$.
The noise satisfies $\vect(\mathbf W)\sim\CN(\mathbf0,v\I_D)$
independently of $\mathbf H$, and both classes use the same sampling
distribution of $v$.
We train the classifier with parameters $\phi$ by minimizing the
binary cross-entropy loss
\begin{equation}
 \mathcal L_{\rm act}(\phi)=
 \E[\log(1+e^{z_\phi})-a z_\phi],
 \label{eq:classifier-loss}
\end{equation}
where $z_\phi=z_\phi(\mathbf R,v,\mathbf F_\theta(\mathbf R,v))$.
For the balanced training distribution, the Bayes-optimal logit is
\begin{equation}
 z^\star(\mathbf R,v)=
       \log\frac{q_{1,v}(\mathbf R)}{q_{0,v}(\mathbf R)}.
 \label{eq:ideal-logit}
\end{equation}

Given a trained classifier $z_\phi$, we use
$\ell_k=z_\phi(\mathbf R_k,v,\mathbf F_k)$ as an estimate of the log
likelihood ratio, where $\mathbf F_k=\mathbf F_\theta(\mathbf R_k,v)$.
Substituting this estimate into \eqref{eq:activity-log-odds} gives the
activity update
\begin{equation}
 p_k=\sigmoid(\logit\rho+\ell_k),
 \label{eq:activity-probability}
\end{equation}
where $\logit a=\log[a/(1-a)]$ and $\sigmoid b=(1+e^{-b})^{-1}$.
